% !TeX spellcheck = it_IT
\newpage
\section{User Experience}
\subsection{Wireframing}
\begin{definition}[Wireframe]
	I wireframe sono \textbf{blueprint} che migliorano la comunicazione tra designer e programmatori riguardo la struttura del sito web o app che stanno sviluppando.
\end{definition}
Servono ad evitare che le varie persone coinvolte nel progetto prendano strade diverse seguendo solo ciò che pensano sia meglio individualmente.\\

L'attività di \textbf{wireframing} è quindi il processo di creazione di sketch basilari dell'interfaccia di un sito o un'applicazione per dimostrare la sua \textbf{struttura} e il suo \textbf{layout}.\\

\noindent Esistono due tipi di wireframe:
\begin{itemize}
	\item \textbf{Low-fidelity}: semplici sketch dell'interfaccia che si concentrano su struttura e layout
	\item \textbf{High-fidelity}: rappresentazione dettagliata dell'interfaccia con topografia di base, icone e dettagli di documentazione
\end{itemize}

\subsection{Adaptability}
Il contesto gioca un ruolo fondamentale nel UX Design (e.g. utente semplice da telefono contro esperto da computer) ed è quindi importante fare \textbf{ricerche sugli utenti} e \textbf{testing} in modo da comprendere a fondo i vari modi in cui gli utenti possono interagire.

\begin{definition}[Adaptability]
	Un termine ombrello nel UX Design che si riferisce alla capacità di un design di adattarsi e rispondere alle diverse necessità, contesti, capacità e preferenze degli utenti.
\end{definition}

\subsubsection{Dispositivo}
\paragraph{Responsive Design} Con sempre più dispositivi venduti, ognuno con diversi layout e standard, è fondamentale fare in modo che i siti web siano leggibili ed utilizzabili ovunque. Il responsive design è un metodo di designing e programmazione di pagine web e app per \textbf{adattarle} a \textit{dimensione} e \textit{risoluzione} dei diversi dispositivi.\\
Il responsive design permette all'utente di avere un'esperienza \textbf{consistente} indipendentemente dal dispositivo e rende più facile \textbf{manutenere} e \textbf{aggiornare} i siti web, dato che va fatto su una sola versione. Esistono alcune tecniche:
\begin{itemize}
	\item \textbf{Flexible grid}: la pagina è divisa in colonne di larghezza uguale e ogni elemento può usarne una o più. Si possono impostare dei breakpoint e degli spazi. Si divide in diversi \textit{pattern}:
	\begin{itemize}
		\item \textbf{Tiny Tweaks}: il layout rimane pressoché lo stresso su tutti i dispositivi e lavora su una colonna, cambiando solo dimensione del font, spaziature e margini
		\item \textbf{Column Drop}: si parte dai dispositivi mobili con una singola colonna e se ne aggiungono man mano che lo schermo cresce, mentre scendono sotto alla precedente per quelli piccoli
		\item \textbf{Mostly Fluid}: la griglia diventa fluida, con i margini che aumentano con il crescere dello schermo e con una larghezza massima per evitare che le righe diventino troppo lunghe
		\item \textbf{Layout Shifter}: il layout cambia significativamente ai breakpoints
		\item \textbf{Off Canvas}: la navigazione o il contenuto secondario vengono nascosti di default sui dispositivi mobili e si possono mostrare tramite menu o gesture
	\end{itemize}
	\item \textbf{Flexible images}: le immagini scalano e si aggiustano in base alla dimensione dello schermo
	\item \textbf{Media query} e \textbf{breakpoints}: permettono l'applicazione di stili diversi in base alle caratteristiche del dispositivo
\end{itemize}
Alcune buone pratiche sono:
\begin{itemize}
	\item \textbf{Keep it simple}: evitare layout complessi e mantenere un design minimale
	\item Priorizzare il \textbf{contenuto}: deve essere la parte più importante della pagina e facilmente leggibile da tutti i dispositivi
	\item Design per il \textbf{touch}: assicurarsi che i bottoni siano grandi abbastanza per essere toccati con un dito
	\item Ottimizzare per il \textbf{tempo di caricamento}, in particolare immagini o altri oggetti pesanti
	\item \textbf{Test} su tanti tipi di dispositivi per verificare che tutto appaia e funzioni correttamente
\end{itemize}

\subparagraph{Mobile first} È una filosofia che priorizza i dispositivi mobili rispetto a quelli desktop o comunque con schermi più grandi. Questo implica che i designer si assicurano che il contenuto e le funzionalità più importanti siano facilmente accessibili sui piccoli schermi e che sia facile navigare il sito con controlli di tipo touch. Inoltre questo approccio si assicura che il sito carichi velocemente sui dispositivi con una connessione più lenta.\\
Si basa su un design di tipo \textbf{progress advancement}, ovvero che parte dalle funzionalità minime e interazioni basilari (per mobile) per poi aggiungerne di più complesse (per desktop). Si oppone alla \textbf{graceful degradation}.

\paragraph{Adaptive Design}
Il design adattivo consiste nel creare diversi tipi di layout per ogni diverso dispositivo. È più costoso del \textit{responsive design} ma a volte è migliore, e.g. grandi aziende che non vogliono riscrivere il codice già presente per desktop.

\subsubsection{Utente}
\paragraph{Adaptive UX} Per far si che un sistema si adatti all'utente è importante avere informazioni su di lui, ad esempio:
\begin{itemize}
	\item \textbf{Collaborative filtering}, e.g. se ad A piace X e Y e a B piace X allora a B potrebbe piacere Y
	\item \textbf{Content-based filtering} che si concentra invece su contenuti simili a qualcosa che piace ad un utente
\end{itemize}
\begin{center}
	\includegraphics[scale=.4]{adaptive_ux}
\end{center}

\paragraph{Accessibilità} È un aspetto fondamentale del UX Design perché assicura che tutti gli utenti possano accedere e interagire con i prodotti digitali indipendentemente dalle loro abilità. Non solo il $15\%$ della popolazione mondiale ha una qualche forma di disabilità, ma bisogna anche tenere in considerazione le compromissioni temporanee (e.g. braccio rotto). In diversi paesi è obbligatorio per legge (e.g. EU).\\
Lo standard è il \textbf{Web Content Accessibility Guidelines}, ad oggi alla versione 2.2, che si compone di quattro principi (POUR):
\begin{itemize}
	\item \textbf{Perceivable}: le informazioni devono essere presentate in modi che l'utente possa percepire
	\item \textbf{Operable}: i componenti dell'interfaccia devono essere utilizzabili
	\item \textbf{Understandable}: le informazioni e le operazioni devono essere chiare
	\item \textbf{Roubst}: il contenuto deve funzionare con tecnologie di supporto
\end{itemize}
Ognuno è poi diviso in tredici linee guida, ognuna con i suoi criteri, che portano a tre livelli di conformità:
\begin{itemize}
	\item \textbf{A}: livello minimo, se non raggiunto alcuni utenti sono completamente bloccati
	\item \textbf{AA}: quello raccomandato e spesso obbligatorio per legge
	\item \textbf{AAA}: il migliore, non sempre raggiungibile a seconda del contenuto
\end{itemize}
Le principali linee guida per il livello \textit{AA} sono:
\begin{itemize}
	\item \textit{Contrasto} $4.5:1$ per il testo e $3:1$ per i componenti dell'interfaccia
	\item Accesso tramite \textit{tastiera} per tutte le funzionalità
	\item \textit{Testo alternativo}, ovvero descrizioni significative per tutte le immagini che non siano decorative
	\item \textit{Semantica HTML} che usi gli heading appropriati, le label per i form di input e ARIA quando necessario
	\item \textit{Consistenza}: navigazione nello stesso ordine, funzionalità uguali chiamate in maniera identica
	\item \textit{Testo ridimensionabile} fino al $200\%$ senza perdere funzionalità
\end{itemize}

\subsection{Wireflows}
A differenza dei \textit{wireframe}, che sono statici, gli \textbf{user flow} aiutano a descrivere le possibili \textbf{azioni} che un utente può fare nell'interfaccia. Aiuta a mappare tutti i possibili passaggi e movimenti dell'utente nell'applicazione o sito web.\\
Partendo dalle \textit{personas}, i \textit{requisiti} e gli \textit{scenari} si pensa a tutti i possibili punti di contatto che l'utente avrà con il sistema per definire così un \textbf{diagramma di flusso}. Questo può essere poi integrato con i \textit{wireframe} per dare vita ai \textbf{wireflows}.

\subsection{Mockups}
\begin{definition}[Mockup]
	Un mockup è una rappresentazione ad alta fedeltà del prodotto che include più dettagli visivi e elementi di design di un wireframe, permettedo di testare e rifinire come il prodotto si presenta.
\end{definition}
I mockup vengono di solito utilizzati per comunicare i concetti di design e le funzionalità in maniera più efficace. Sono anche utili per alcuni tipi di studi sull'usabilità dove serve il feedback sull'aspetto generale del design.

\begin{note}
	È fondamentale notare che un mockup NON ha le funzionalità necessarie per fare test significativi od ottenere feedback profondi in quanto non possiede alcun elemento interattivo.
\end{note}

\subsection{Pretotyping}
\subsubsection{Mercato}
Così come in legge un imputato e innocente fino a prova contraria, nel mercato dovremmo assumere che un prodotto sia fallito fino a prova contraria. Questo perché statisticamente solo il $10\%$ dei prodotti ha successo.
\begin{definition}[Market failure]
	La maggior parte dei nuovi prodotti fallirà sul mercato, anche se sono eseguiti in maniera competente.
\end{definition}
Il pretotyping fornisce strumenti e tecniche necessarie a validare un'idea con poche risorse e in poco tempo prima di investire per lo sviluppo.\\

Per minimizzare il rischio di falsi negativi o positivi sul mercato è necessario partire da "\textbf{thoughtland}" per ottenere opinioni
\begin{observation}[Translation problem]
	Un'idea è un'astrazione soggettiva che ci si immagina nella propria testa. Nel momento in cui la si prova a comunicare c'è un problema di "traduzione", sopratutto se è molto innovativa.
\end{observation}
\begin{observation}[Prediction problem]
	Anche se gli interlocutori hanno un'idea astratta della tua idea che è molto simile a come la immagini, gli umani sono pessimi a predire se vogliono veramente qualcosa che non hanno mai sperimentato e come la userebbero.
\end{observation}
per poi arrivare ad "\textbf{actionland}" per ottenere dati.

\subsubsection{Pretotype}
Prima ancora di creare un \textit{prototipo} è importante capire se fallirà o meno sul mercato prima di investire tempo e denaro. Questo si fa con un \textbf{pretotype}. I passaggi per lo sviluppo sono:
\begin{enumerate}
	\item Isolare l'\textbf{assunzione chiave} e capire se è vera o falsa
	\item Scegliere un tipo di \textbf{pretotype} che permetta di isolare l'assunzione
	\item Fare un'ipotesi di \textbf{interesse di mercato}: quante persone e di che tipo vorrebbero il pretotype
	\item \textbf{Test} del pretotype mettendolo nel mondo reale e analizzando come le persone ci interagiscono
	\item \textbf{Imparare} dai risultati del test per \textbf{rifinire} il pretotype. Se l'ipotesi permane, testare in nuove situazioni (\textbf{hypozoom})
\end{enumerate}

\paragraph{Fake Door} Questo metodo consiste nel pubblicizzare l'idea del prodotto che ancora non esiste fingendo il contrario e vedere quante persone sarebbero interessate. Si può fare tramite pubblicità, siti web, form, email o poster.

\paragraph{Mechanical Turk} Fingere che un algoritmo sia implementato facendo in realtà fare il lavoro ad un umano e vedere se interessa.

\paragraph{Impersonator} Utilizzare un prodotto già esistente con un'interfaccia diversa.

\paragraph{Pinocchio} Creare una finta versione del prodotto che non funziona per verificare ad esempio se è il giusto form factor, ha le giuste funzionalità ed è usabile. E.g. finto telefono di plastica.

\paragraph{One Night Stand} Fare una versione completa dell'esperienza che si vuole creare ma solo per una volta per capire se può funzionare.

\paragraph{Facade} Noleggiare o chiedere in prestito attrezzature, spazi e materiali per simulare un'infrastruttura che in realtà non esiste, evitando così di fare subito un investimento grosso.

\subsection{Prototyping}
Nel momento in cui si ha imparato tutto il possibile sulla domanda di mercato dai \textit{pretotypes} ed è necessario mettere il vero prodotto nelle mani dei clienti per testarlo e rifinirlo, si crea un \textbf{prototipo}.

\begin{definition}[Prototipo]
	Un primo esempio sperimentale di un design che permette all'utente di visualizzare e \textbf{interagirci} prima che il prodotto finale sia sviluppato.
\end{definition}

\begin{definition}[Prototyping]
	Il processo con il quale i team di design ideano, esperimentano e portano in vita i concetti. È la fase in cui le idee sono implementate in un qualcosa di tangibile.
\end{definition}

Un prototipo porta quindi a diversi vantaggi:
\begin{itemize}
	\item Fornisce delle fondamenta solide per l'ideazione e dà agli stakeholder una chiara immagine dei potenziali benefici, rischi e costi
	\item I cambiamenti possono essere fatti all'inizio del processo, prima del rilascio, evitando errori costosi e il dover rimanere per forza attaccati ad una versione
	\item Il feedback dell'utente permette di identificare cosa funziona meglio e se è necessario rifare delle parti
\end{itemize}

\subsubsection{Fidelity}
\begin{definition}[Fidelity]
	Quanto un prototipo assomiglia al prodotto finale in termini di aspetto e funzionalità.
\end{definition}
Esistono due livelli di fedeltà di un prototipo:
\begin{itemize}
	\item \textbf{Low-fidelity}: di solito molto economici in termini di soldi e tempo, fornisce interazioni molto semplici (se non inesistenti). Facile da modificare. È poco realistico e potrebbe essere troppo basilare perché un utente riesca a dare un feedback, senza contare che li costringe a dover immaginare come sarebbe il prodotto.
	\item \textbf{High-fidelity}: più costosi e lunghi da creare, sono spesso sviluppati in ambienti specifici (e.g. Figma) e danno una sensazione di prodotto finito, rendendoli attraenti sia per gli utenti che per gli stakeholder. I test forniscono dati più accurati anche se gli utenti potrebbero soffermarsi troppo su dettagli superficiali invece che il contenuto.
\end{itemize}
Ed eventualmente tutte le possibilità che ci sono in mezzo (\textbf{mid-fidelity}), ad esempio un wireframe avanzato con solo alcune parti interattive. In generale la fedeltà si basa su sei dimensioni:
\begin{itemize}
	\item \textbf{Realismo fisico/visuale}: quanto simile è il prototipo al prodotto finale per quanto riguarda l'estetica e il design visivo
	\item \textbf{Obiettivo}: l'ampiezza e profondità del design rappresentato nel prototipo
	\item \textbf{Funzionalità}: cosa funziona effettivamente nel prototipo, e.g. bottoni in una web app
	\item \textbf{Dati}: se il prototipo funziona con dati reali o finti
	\item \textbf{Autonomia}: se il prototipo funziona da solo o ha bisogno di input dal designer o dall'utente
	\item \textbf{Piattaforma}: se il prototipo è un interim o l'implementazione finale
\end{itemize}

\subsubsection{Scope}
Lo \textit{scope} descrive quanto l'utente sarà in grado di vedere e testare. Si divide in:
\begin{itemize}
	\item \textbf{Orizzontale}: fornisce una visione completa del sistema senza però andare troppo nel dettaglio per ogni funzionalità. È quindi limitato ad un solo layer. È utile agli stadi iniziali del design in quanto permette agli stakeholder di identificare velocemente potenziali problemi.
	\item \textbf{Verticale}: fornisce una visione dettagliata di una particolare specifica, permettendo di testare se il prodotto è fattibile o se un particolare requisito è soddisfatto
\end{itemize}

\begin{center}
	\includegraphics[scale=.25]{scope}
\end{center}

\subsubsection{Tecniche}
\paragraph{Carta} Un prototipo di carta è \textit{low-fidelity} e rappresenta il layout e le funzionalità basilari di un prodotto o servizio. È veloce ed economico. Consiste nel disegnare su carta o cartone e poi l'utente può interagirci fisicamente. Può essere usato per testare agli stadi iniziali dove posizionare bottoni e menu per il flusso delle informazioni e le interazioni dell'utente.

\paragraph{Wireframe} È la versione digitale di quello di carta.

\paragraph{Mechanical Turk} Un prototipo che simula il comportamento del prodotto usando in realtà input umani invece della tecnologia. Ad esempio può consistere in audio o video preventivamente registrato, risposte programmate o interazioni live con l'utente.

\paragraph{Funzionali} Prototipo  che include funzionalità funzionanti o parzialmente funzionanti per testare gli aspetti tecnici di un prodotto, come reattività o funzionamento. Utile negli stadi finali del design quando si vogliono perfezionare dettagli tecnici prima di passare allo sviluppo completo.

\paragraph{Non funzionali} Un prototipo che sembra esattamente come il prodotto reale ma non ha nessuna funzionalità. È utilizzato per dimostrare il design fisico e l'interfaccia utente senza dover programmare in maniera dettagliata.

\begin{observation}
	La differenza tra i prototipi e i \textit{mockup} e \textit{wireframe} è che i primi sono versioni più avanzate creati sulla base degli ultimi che permettono di testare e validare il design prima della fase di sviluppo.
	\begin{center}
		\includegraphics[scale=.3]{confronto}
	\end{center}
\end{observation}

\subsection{Design systems}
Con la crescita di un prodotto più designer prendono decisioni diverse e i componenti vengono reimplementati in maniera leggermente differente. Un po' alla volta le \textbf{inconsistenze} si accumulano e la qualità scende.

\begin{definition}[Design System]
	Un insieme di componenti riutilizzabili, guidati da standard precisi, che possono essere assemblati per costruire applicazioni. Include:
	\begin{itemize}
		\item Libreria dei componenti, e.g. bottoni
		\item Design tokens, e.g. colori e font
		\item Linee guida per l'utilizzo e documentazione
		\item Principi di design e pattern
		\item Implementazioni di codice
	\end{itemize}
\end{definition}

\noindent I sistemi di design garantiscono:
\begin{itemize}
	\item \textbf{Consistenza}: gli stessi componenti si comportano allo stesso modo ovunque
	\item \textbf{Efficienza}: crea una volta e riusa molteplici
	\item \textbf{Qualità}: accessibilità e responsività
	\item \textbf{Scalabilità}: nuovi membri del team si possono unire più facilmente
	\item \textbf{Manutenibilità}: risolvere un errore lo fa per tutti
\end{itemize}

\subsubsection{Atomic Design}
Questo metodo si compone di blocchi da costruzione uno più grande dell'altro:
\begin{itemize}
	\item \textbf{Atomi}: elemento basilare che non può essere diviso ulteriormente, e.g. bottoni
	\item \textbf{Molecole}: combinazioni di atomi che funzionano assieme, e.g. form di ricerca (bottone, input e label)
	\item \textbf{Organismo}: componenti complessi dell'interfaccia composti dai precedenti, e.g. header
	\item \textbf{Template}: layout di pagina che mostrano la struttura del contenuto, e.g. homepage
	\item \textbf{Pagine}: istanze specifiche di layout con vero contenuto
\end{itemize}

\subsubsection{Design Tokens}
\begin{definition}[Design Token]
	Decisione di design salvata come dato.
\end{definition}
Permette di cambiare una volta per tutte le occorrenze e indipendentemente dalla piattaforma. E.g. definire il colore primario.

\subsection{Evaluation}
Una volta creato un \textit{prototipo} lo si vuole \textbf{testare} con gli utenti. Questo può portare a trovare problemi di \textbf{usabilità}, possibile margine di miglioramento e capire come gli utenti si comportano e cosa preferiscono.

\subsubsection{Definizioni}
Esistono due definizioni diverse di \textit{usability} che si concentrano su due aspetti leggermente diversi.
\paragraph{Nielsen}
\begin{definition}[Usability - Nielsen]
	L'usabilità è un attributo di qualità che definisce quanto facili le interfacce utente sono da utilizzare. Si compone di:
	\begin{itemize}
		\item \textbf{Learnability}: quanto è facile compiere azioni basilari per la prima volta
		\item \textbf{Efficiency}: una volta imparato il design, quanto veloce si eseguono le azioni
		\item \textbf{Memorability}: quando gli utenti ritornano sul prodotto dopo tanto tempo, quanto è facile riprenderne l'utilizzo
		\item \textbf{Errors}: quanti errori compie l'utente, quanto sono gravi e quanto è facile riprendersi da essi
		\item \textbf{Satisfaction}: quanto è piacevole usare il design
	\end{itemize}
\end{definition}

\paragraph{ISO} La definizione ISO si concentra invece sul fatto che l'usabilità sia il risultato dell'interazione con il sistema: i vari attributi possono aiutare ma l'usabilità non è una caratteristica intrinseca del prodotto.
\begin{definition}[Usability - ISO9241-11]
	L'usabilità è il grado a cui un sistema, prodotto o servizio permette agli utenti di compiere determinate azioni con \textbf{efficacia}, \textbf{efficienza} e \textbf{soddisfazione} in uno specifico contesto d'uso.
\end{definition}

\begin{definition}[Efficacia]
	Accuratezza e completezza con le quali gli utenti raggiungono obiettivi specifici. Si può misurare in:
	\begin{itemize}
		\item Percentuale di task completate con successo
		\item Tasso di errore
		\item Accuratezza nelle decisioni, e.g. percentuale di selezioni corrette
		\item Completezza dei risultati
	\end{itemize}
\end{definition}

\begin{definition}[Efficienza]
	Risorse utilizzate in relazione ai risultati ottenuti, e.g. tempo, impegno umano, costi o materiali. Si misura in:
	\begin{itemize}
		\item Mediana o percentile del tempo speso per ogni task completata con successo
		\item Throughput delle task per ora o passi necessari per completarla
		\item Impegno umano valutato tramite scale di workload, e.g. NASA-TLX, SWAT o MRQ
		\item Costo per risultato corretto quando rilevante
	\end{itemize}
\end{definition}

\begin{definition}[Soddisfazione]
	Quanto le risposte fisiche, cognitive ed emotive di un utente raggiungono le sue necessità e aspettative. Si misura in:
	\begin{itemize}
		\item \textbf{System Usability Scale} per l'usabilità generale percepita
		\item \textbf{UMUX-Lite} per l'utilità/usabilità percepita quando \textit{SUS} è troppo lungo
		\item Scale specifiche di usabilità come \textbf{Chatbot Usability Scale}
	\end{itemize}
\end{definition}

\begin{note}
	Conviene separare la \textit{soddisfazione} (soggettiva) dall'\textit{efficacia/usabilità} (oggettiva) ma analizzarle assieme per provare a spiegare le divergenze, e.g. veloce ma frustrante.
\end{note}

\paragraph{Utile}
\begin{definition}[Utilità]
	Quanto un prodotto riesce a fare ciò di cui gli utenti hanno bisogno.
\end{definition}
\noindent La combinazione di \textit{utilità} e \textit{usabilità} india che un prodotto è \textbf{utile}.

\subsubsection{Valutazioni euristiche}
La valutazione euristica è un metodo di valutazione delle interfacce utente basato su principi di usabilità ben definiti (\textit{euristiche}). Di solito comprende un piccolo gruppo di valutatori, circa $3-5$. Anche se questo metodo non fornisce un approccio sistematico alla risoluzione dei problemi di usabilità, può \textbf{guidare} la creazione di design rivisitati basati sui principi violati trovati in maniera economica.

\begin{center}
	\includegraphics[scale=.3]{euristiche}
\end{center}

\noindent Nielsen ha identificato dieci euristiche per l'usabilità:
\begin{enumerate}
	\item \textbf{Visibilità} dello \textbf{stato del sistema}: i design dovrebbero mantenere l'utente informato riguardo a ciò che sta accadendo tramite feedback appropriati
	\item \textbf{Mappare} il sistema al \textbf{mondo reale}: il design dovrebbe usare parole, frasi e concetti familiari per l'utente
	\item \textbf{User control} e \textbf{libertà}: gli utenti necessitano di una "uscita di emergenza" evidente per poter abbandonare l'azione non voluta
	\item \textbf{Consistenza} e \textbf{standard}: seguire le convenzioni della piattaforma per evitare che l'utente debba provare a indovinare il significato delle cose
	\item \textbf{Prevenzione degli errori}: prima ancora di avere dei messaggi di errore ben strutturati è importante provare ad evitare che accadano
	\item \textbf{Recognition} invece di recall: ridurre al minimo il carico sulla memoria degli utenti facendo in modo che gli elementi, le azioni e le opzioni siano visibili
	\item \textbf{Flessibilità} ed \textbf{efficienza} d'uso: garantire shortcut che un utente esperto può sfruttare
	\item \textbf{Estetica} e \textbf{design minimal}: evitare informazioni irrilevanti nelle interfacce
	\item \textbf{Riconoscere}, \textbf{diagnosticare} e \textbf{riprendersi} dagli \textbf{errori}: i messaggi di errore dovrebbero essere frasi che indicano precisamente il problema e suggeriscono una soluzione costruttiva, invece che un codice
	\item \textbf{Aiuto} e \textbf{documentazione}: per quanto idealmente il design non dovrebbe aver bisogno di ulteriori spiegazioni, potrebbe servire della documentazione da fornire all'utente per completare certi task
\end{enumerate}

\subsubsection{Cognitive Walkthrough}
Metodo di valutazione che si base sull'idea che l'utente impari ad usare l'interfaccia tramite l'\textbf{esplorazione}, aiutato dai suoi \textbf{obiettivi} e utilizzando semplice \textbf{ragionamento}.\\

La valutazione viene quindi fatta da un \textbf{esperto} che compie un insieme di task imitando le performance di un utente. Prima che l'esperto possa procedere, deve ricevere dal designer o sviluppatore:
\begin{itemize}
	\item \textbf{Descrizione utente}, che include il suo livello di esperienza con i computer e le assunzioni del designer
	\item \textbf{Descrizione sistema} che include le operazioni e le performance
	\item \textbf{Descrizione task} che specifica ciò che l'esperto dovrà fare seguendo il punto di vista dell'utente
	\item \textbf{Sequenza azioni} che descrive ciò che il sistema deve mostrare e le azioni che l'utente deve compiere per completare una task. La combinazione di un system display e un'azione utente compongono uno \textit{step}
\end{itemize}
Durante ogni step della \textit{sequenza azioni}, l'esperto deve porre le seguenti domande:
\begin{itemize}
	\item A questo stadio, il prossimo obiettivo è chiaro?
	\item L'azione corretta è ovvia?
	\item È chiaro che questa azione porta all'obiettivo?
	\item Quali problemi ci sono nel compiere questa azione?
\end{itemize}

\newpage
\subsubsection{Usability testing}
Processo nel quale un \textit{ricercatore} chiede ad un \textit{partecipante} di compiere delle task, di solito usando una o più interfacce, mentre il suo comportamento viene \textbf{osservato} e ascoltato per eventuali \textbf{feedback}.
\begin{wrapfigure}[17]{r}{7cm}
	\vspace{1cm}
	\includegraphics[width=7cm]{elements_usability}
\end{wrapfigure}
Abbiamo quindi tre elementi:
\begin{itemize}
	\item \textbf{Facilitator}: il ricercatore che guida il \textit{partecipante} durante il test. Deve dare istruzioni, rispondere alle domande e farne in seguito. Il suo lavoro è assicurare che i risultati siano validi e di alta qualità senza influenzare il comportamento del \textit{partecipante}.
	\item \textbf{Task}: attività realistiche che il \textit{partecipante} potrebbe dover compiere nella vita reale; possono essere molto specifiche o molto generali in base alle domande del ricercatore e al tipo di test. Possono essere comunicate verbalmente o in forma scritta ma è importante che siano precise al massimo in quanto il minimo errore può influenzare il \textit{partecipante}.
	\item \textbf{Partecipante}: utente realistico del prodotto o servizio; può essere qualcuno che lo usa già o comunque che ha un simile background o necessità. Idealmente la scelta parte dalle \textit{personas} definite. A volte gli viene chiesto di pensare ad alta voce.
\end{itemize}
Esistono due tipi di test: \textbf{quantitativo} e \textbf{qualitativo}.

\begin{table}[h]
	\centering
	\renewcommand{\arraystretch}{1.5}
	\small
	\begin{tabular}{|P{3.5cm}|P{5.5cm}|P{5.5cm}|}
		\hline
		% Titoli centrati, in grassetto e senza colore di sfondo
		\multicolumn{1}{|c|}{} & 
		\multicolumn{1}{c|}{\textbf{Ricerca Qualitativa}} & 
		\multicolumn{1}{c|}{\textbf{Ricerca Quantitativa}} \\ \hline
		
		\textbf{Domande a cui risponde} & Perché? & Quanti e quanto? \\ \hline
		
		\textbf{Obiettivi} & 
		Sia \textit{formativa} che \textit{riepilogativa}:
		\begin{itemize}[nosep, leftmargin=*]
			\item informare le decisioni di design
			\item identificare problemi di usabilità e trovare soluzioni
		\end{itemize} & 
		Principalmente \textit{riepilogativa}:
		\begin{itemize}[nosep, leftmargin=*]
			\item valutare l'usabilità di un sito esistente
			\item monitorare l'usabilità nel tempo
			\item confrontare il sito con i competitor
			\item calcolare il ROI
		\end{itemize} \\ \hline
		
		\textbf{Quando viene usata} & 
		In qualsiasi momento: durante il redesign, o quando si ha un prodotto finale funzionante. & 
		Quando si ha un prodotto funzionante (all'inizio o alla fine di un ciclo di design). \\ \hline
		
		\textbf{Risultati} & 
		Risultati basati sulle impressioni, interpretazioni e conoscenze pregresse del ricercatore. & 
		Risultati statisticamente significativi che probabilmente saranno replicati in uno studio diverso. \\ \hline
		
		\textbf{Metodologia} & 
		\begin{itemize}[nosep, leftmargin=*]
			\item Pochi partecipanti
			\item Condizioni di studio flessibili adattabili alle esigenze del team
			\item Protocollo "think-aloud"
		\end{itemize} & 
		\begin{itemize}[nosep, leftmargin=*]
			\item Molti partecipanti
			\item Condizioni di studio ben definite e rigorosamente controllate
			\item Solitamente senza "think-aloud"
		\end{itemize} \\ \hline
	\end{tabular}
\end{table}

\newpage
\paragraph{Pianificazione} Prima di tutto bisogna definire l'\textbf{obiettivo} della ricerca, ovvero che cosa vogliamo testare e a che scopo. Successivamente va deciso il formato, che può essere:
\begin{itemize}
	\item In \textbf{laboratorio} o \textbf{sul campo}
	\item \textbf{Moderato} o \textbf{non moderato}: il primo caso fornisce maggiori informazioni e un miglior accesso ai commenti ma costa di più ed è più difficile da organizzare. Il secondo è più economico, veloce e fornisce un miglio accesso a partecipanti difficili da reclutare
	\item \textbf{Di persona} o da \textbf{remoto}: il primo è il più consigliato quando possibile in quanto fornisce un'interazione più personale che aiuta a capire indizi sottili
\end{itemize}
Per quanto riguarda il \textbf{numero di partecipanti}, Nielsen evidenzia che per studi \textit{qualitativi} $5$ persone sono un buon numero. Per quanto riguarda quelli \textit{quantitativi} dovrebbero essere almeno $20$. In entrambi i casi devono essere i partecipanti corretti che rispecchiano la demografia e i comportamenti dei \textit{target user}.
\begin{note}
	Evita di chiedere ai partecipanti di fingere o immaginare scenari in quanto potrebbe invalidare i risultati.
\end{note}
\noindent Il passo successivo è quello di scrivere le giuste \textbf{task}, che possono essere:
\begin{itemize}
	\item \textbf{Esplorative}, utilizzate per imparare come le persone esplorano le informazioni; non adatte alla ricerca \textit{quantitativa}
	\item \textbf{Specifiche}: sono mirate e hanno una risposta corretta o un punto di arrivo
\end{itemize}
Infine si decidono le \textbf{metriche}, se presenti. Nel caso degli studi \textit{qualitativi} si definiscono le domande da fare ai partecipanti. In quelli \textit{quantitativi} sono invece fondamentali, e.g. tempo speso su una task, livelli di soddisfazione, ratio di successo e di errore.

\begin{note}
	Può essere molto utile fare prima un \textbf{pilot study} per rifinire come sono scritte le task, il loro ordine e quantità e assicurarsi che si stiano reclutando i giusti partecipanti.
\end{note}

\paragraph{Cosa evitare} Ecco dieci cose da evitare durante il testing:
\begin{enumerate}
	\item Dire agli utenti \textbf{dove muoversi} nell'interfaccia
	\item Dire agli utenti \textbf{cosa fare}
	\item Creare task con \textbf{date passate}
	\item Rendere le task \textbf{troppo semplici}
	\item Usare scenari \textbf{troppo elaborati}
	\item Scrivere una \textbf{pubblicità} invece di un task
	\item Rischiare di causare una \textbf{reazione emotiva}
	\item Essere \textbf{divertenti}
	\item \textbf{Offendere} i partecipanti
	\item \textbf{Chiedere come} un partecipante farebbe qualcosa
\end{enumerate}

\newpage
\subsection{IOT}
\begin{definition}[Internet of Things]
	Un sistema di \textit{computing devices} interconnessi, macchine meccaniche e digitali, oggetti, animali o persone che sono dotati di \textbf{unique identifiers} (UIDs) e l'abilità di trasferire dati attraverso una rete senza la necessità di interazione umana.
\end{definition}
\begin{definition}[Industria 4.0]
	Descrive l'organizzazione del processo di produzione basato sulla comunicazione automatica di tecnologia e dispositivi lungo la catena produttiva.
\end{definition}
\begin{center}
	\includegraphics[scale=.4]{industria4}
\end{center}
Nell'industria 4.0 i sistemi monitorano i processi finisci e creano una copia virtuale del mondo fisico per poi prendere decisioni decentralizzate basate su meccanismi di auto organizzazione.
\begin{definition}[Digital Twin]
	Una copia digitale di un'entità fisica come processi, persone, luoghi, sistemi e dispositivi.
\end{definition}
\subsubsection{Il Ciclo di Vita dell'Esperienza}
Il modello analizzato suddivide l'interazione con i sistemi IoT in tre fasi temporali, che superano il semplice momento dell'acquisto:
\begin{itemize}
	\item \textbf{Before Use (Anticipated use)}: la fase di anticipazione e aspettativa verso le potenzialità del dispositivo.
	\item \textbf{During Use (Actual use)}: l'interazione tecnica diretta, dove l'IoT deve garantire efficacia e soddisfazione nel completamento dei task.
	\item \textbf{After Use (Digested use)}: la fase di elaborazione mentale e ritenzione emotiva, cruciale per la fidelizzazione in ecosistemi di servizi ricorsivi.
\end{itemize}

\subsubsection{Livelli di Design e Integrazione Aziendale}
La progettazione per l'IoT non si limita all'interfaccia, ma si estende a livelli strategici più profondi:
\begin{description}
	\item[Usability]: Focalizzata sull'efficienza del compito durante l'uso effettivo del dispositivo.
	\item[User Experience (UX)]: La progettazione dell'esperienza globale legata al singolo prodotto connesso, coprendo l'intero arco temporale (prima, durante e dopo).
	\item[Customer Experience (CX)]: La gestione di tutti i \textit{touchpoint} dell'ecosistema (es. Advertising, Customer Service, Ethics), vedendo il dispositivo IoT come uno dei molti canali di comunicazione azienda-cliente.
	\item[Service Design]: Il livello più alto di astrazione che coordina la strategia di business, il personale e il design organizzativo per sostenere l'erogazione del servizio nel tempo.
\end{description}

In sintesi, l'IoT trasforma il prodotto in un'interfaccia di servizio permanente, rendendo il \textit{Service Design} il pilastro fondamentale per la coerenza tra la promessa tecnologica e l'esperienza realmente vissuta dall'utente.

\subsubsection{User Experience}
Il design per l'IoT è per natura più \textbf{complesso} di quello per i servizi web. Questo a causa del livello attuale della tecnologia, dell'immatura comprensione del valore percepito dagli utenti e dal semplice fatto che ci sono più aspetti da considerare, come:
\begin{itemize}
	\item La natura \textbf{specializzata} dei dispositivi IoT: molti hanno bisogno di una forma fisica precisamente progettata e ingegnerizzata e di conseguenza necessitano di più funzionalità di UI
	\item L'abilità dei dispositivi di connettere il mondo digitale con quello fisico
	\item Il fatto che molti prodotti sono \textbf{sistemi distribuiti} composti da più dispositivi: bisogna capire come distribuire le funzionalità e progettare UI e interazioni su tutto il sistema come un'unica entità, garantendo \textit{coerenza} (\textbf{interusability})
	\item Il \textbf{networking}: nei dispositivi IoT non è sempre detto che ci sia connessione, cosa che invece era assunta nei servizi web
	\item \textbf{Remote configuration}: se prima l'utente aveva una rappresentazione visuale e immediata degli oggetti, con i dispositivi IoT no
	\item \textbf{Power saving}
\end{itemize}

\subsection{Artificial Intelligence}
L'intelligenza artificiale è sempre più diffusa in tutti gli ambiti. Il problema principale è che, che funzioni o meno, gli utenti \underline{non capiscono} come o perché. Richiede quindi degli approcci fondamentalmente diversi di User Experience.

\subsubsection{Differenze}
A differenza del software tradizionale, la IA è \textbf{probabilistica}, fornisce \textbf{falsi positivi} e \textbf{negativi}, si basa su un principio \textbf{black-box} dove il processo decisionale è nascosto e ha \textbf{comportamenti imprevedibili} che possono essere confusionari, pericolosi o distruttivi.

\begin{definition}[Human-Centered AI]
	AI  che amplifica, aumenta e migliora le performance umane mantenendo i sistemi affidabili e sicuri.
\end{definition}
Questo approccio vuole quindi far \textbf{collaborare} umani e IA in maniera trasparente e con livelli di controllo appropriati. Inoltre è fondamentale lasciare alcune sezioni dell'interfaccia \textit{deterministiche} in modo che, anche quando l'IA non è sicura (che dovrebbe essere il meno possibile), gli utenti possono ricadere su interazioni prevedibili.

\subsubsection{Affidamento}
Il problema dell'affidamento è che gli utenti potrebbero usare l'IA \textbf{troppo} o \textbf{ignorarla completamente} anche se potrebbe essere utile. È fondamentale trovare la giusta via di mezzo basata anche sul livello di \textit{accuratezza} e sull'importanza delle conseguenze.\\

A differenza della \textbf{fiducia}, che è psicologica  e difficile da misurare, l'\textit{affidamento} è osservabile, quantificabile e può essere calibrato in base alle performance.

Per costruire un affidamento appropriato bisogna fornire questi tre meccanismi:
\begin{itemize}
	\item \textbf{Explainability}: aiutare gli utenti a capire il perché mostrando ciò che un sistema può fare e come gli input influenza i risultati, rimanendo specifici sui livelli di performance
	\item \textbf{Trasparenza}: aiutare gli utenti a capire quando affidarsi, mostrando le performance per categoria, i livelli di confidence delle predizioni e avvisando quando la situazione è diversa dai dati di training
	\item \textbf{Controllo}: gli utenti devono poter controllare i dati di input, le informazioni personali e i risultati
\end{itemize}
È importante capire che \textbf{augmentation} (AI per suggerire alternative ed eseguirle se approvate) e \textbf{automation} (AI prima esegue e poi informa quando vuole) esistono su uno \textbf{spettro}. Bisogna capire a che livello dare controllo agli umani.
\begin{center}
	\includegraphics[scale=.3]{automation}
\end{center}

\subsubsection{Miglioramento}
L'IA impara ed evolve nel tempo consumando dati e ricevendo \textbf{feedback} sugli output, che possono essere di tre tipi feedback:
\begin{itemize}
	\item \textbf{Esplicito}: l'utente lo dice esplicitamente per migliorare il sistema, e.g. like/dislike
	\item \textbf{Implicito}: dedotto dal comportamento dell'utente e le sue azioni, e.g. skip rate
	\item \textbf{Duale}: combinazione dei due
\end{itemize}
È importante fare la scelta in base al contesto, impegno e qualità del segnale necessari.\\

Nella pratica il feedback deve essere semplice da dare, mostrare subito l'impatto che ha avuto e rispettare il contesto, quindi non essere utilizzato troppo.

\subsubsection{Design}
Vediamo tre principi fondamentali:
\begin{itemize}
	\item \textbf{Aspettative chiare}: mostrare cosa il sistema può e non può fare, con quale accuratezza e limitazioni e come fa ad imparare
	\begin{note}
		Sempre meglio dare basse aspettative e poi fare più del previsto, non il contrario.
	\end{note}
	\item \textbf{Design per gli errori}: esistono diversi tipi di errore
	\begin{itemize}
		\item \textbf{System}: problemi tecnici e crash
		\item \textbf{Model}: predizioni sbagliate e allucinazioni
		\item \textbf{Data}: dati di training di bassa qualità o input prevenuti
		\item \textbf{Relevance}: risultato corretto ma irrilevante
		\item \textbf{User}: input non chiaro, errore di comprensione
	\end{itemize}
	Per gestirli bisogna indicare quando si presentano, spiegare cosa è andato storto e offrire opportunità per correggere. Bisogna sempre permettere di tornare indietro e  dare il controllo all'utente.
	\item \textbf{Bilanciare complessità} e \textbf{controllo}: mostrare prima le funzionalità base e aggiungere le altre nel tempo, sempre rispettando l'autonomia dell'utente e la possibilità di scartare le decisioni dell'IA
\end{itemize}

\newpage
\noindent I requisiti di design cambiano molto anche in base a quanto c'è in gioco:
\begin{itemize}
	\item \textbf{High-stakes}: serve massima trasparenza, controllo ed explainability in quanto un errore potrebbe avere conseguenze catastrofiche (e.g. medicina)
	\item \textbf{Medium-stakes}:  serve una buona trasparenza e moderato controllo, e.g. educazione
	\item \textbf{Low-stakes}: serve un livello base di trasparenza, e.g. autocompletamento
\end{itemize}

\paragraph{Precision e Recall}
\begin{table}[h]
	\centering
	\renewcommand{\arraystretch}{2}
	\begin{tabular}{|l|c|c|}
		\hline
		& \textbf{Positivo Predetto} & \textbf{Negativo Predetto} \\ \hline
		\textbf{Positivo Reale} & 
		\begin{tabular}[c]{@{}c@{}} \textbf{Vero Positivo (VP)} \\ Predizione corretta \end{tabular} & 
		\begin{tabular}[c]{@{}c@{}} \textbf{Falso Negativo (FN)} \\ Rilevamento mancato \end{tabular} \\ \hline
		\textbf{Negativo Reale} & 
		\begin{tabular}[c]{@{}c@{}} \textbf{Falso Positivo (FP)} \\ Falso allarme \end{tabular} & 
		\begin{tabular}[c]{@{}c@{}} \textbf{Vero Negativo (VN)} \\ Rifiuto corretto \end{tabular} \\ \hline
	\end{tabular}
\end{table}
\begin{definition}[Precisione]
	Quando ha trovato qualcosa, quanto spesso è corretto?
	\begin{equation*}
		\frac{\text{True Positive}}{\text{True Positive} + \text{False Positive}}
	\end{equation*}
\end{definition}

\begin{definition}[Recall]
	Di tutte le cose che dovrebbe trovare, quante ne ha trovate effettivamente?
	\begin{equation*}
		\frac{\text{True Positive}}{\text{True Positive} + \text{False Negative}}
	\end{equation*}
\end{definition}

È importante ottimizzare uno dei due parametri a seconda del \textbf{contesto}: la \textit{precisione} serve quando la cosa peggiore è un falso allarme (e.g. spam detecion) mentre la \textit{recall} quando è un positivo mancante (e.g. analisi mediche).

\paragraph{Processo} Il processo di design per la IA itera su tre cose contemporaneamente: UI, modello di IA e dati. Di seguito i passaggi:
\begin{enumerate}
	\item \textbf{Comprendere} le \textbf{necessità} dell'utente, prestando attenzione al \textit{livello di controllo}, quanto c'è in gioco e come l'utente costruirà la fiducia nel tempo
	\item \textbf{Prototype} anticipato con un test di tipo \textit{Wizard of Oz} per valutare i pattern di interazione
	\item \textbf{Definire i requisiti} del modello collaborando con data scientists e misurando varie metriche tra cui anche le performance umane e la soddisfazione degli utenti
	\item \textbf{Test} con dati veri: la qualità dei dati è fondamentale e bisogna trovare strategie per ottenerli
	\item \textbf{Collaborazione} durante tutto il processo tra gli sviluppatori e i designer
\end{enumerate}